\documentclass[10pt]{article}

% ============================================================
% ECONOMICS LATEX PREAMBLE & SETUP QUICK LOOKUP
% The PDF intentionally PRINTS setup code for easy reference.
% ============================================================

% ---------- Core packages used to typeset THIS handbook ----------
\usepackage[margin=0.68in]{geometry}
\usepackage{amsmath,amssymb,amsfonts,mathtools,bm}
\usepackage{booktabs,array,tabularx,longtable}
\usepackage{graphicx}
\usepackage{amsthm}
\usepackage{siunitx}
\usepackage{threeparttable}
\usepackage{xcolor}
\usepackage{listings}
\usepackage{enumitem}
\usepackage{microtype}
\usepackage[hidelinks]{hyperref}
\usepackage{titlesec}
\usepackage{fancyhdr}

% ---------- Handbook typography ----------
\setlength{\parindent}{0pt}
\setlength{\parskip}{4pt}
\setlist{nosep,leftmargin=*}
\setlength{\tabcolsep}{5pt}
\renewcommand{\arraystretch}{1.08}
\titleformat{\section}{\large\bfseries}{\thesection.}{0.55em}{}
\titleformat{\subsection}{\normalsize\bfseries}{\thesubsection}{0.55em}{}
\titlespacing*{\section}{0pt}{8pt}{4pt}
\titlespacing*{\subsection}{0pt}{6pt}{2pt}

\pagestyle{fancy}
\fancyhf{}
\lhead{LaTeX Quick Lookup for PhD Students}
\rhead{Preamble \& Setup}
\cfoot{\thepage}
\renewcommand{\headrulewidth}{0.3pt}

% ---------- Code appearance ----------
\definecolor{codegray}{RGB}{246,246,246}
\definecolor{rulegray}{RGB}{210,210,210}
\lstset{
  language=[LaTeX]TeX,
  basicstyle=\ttfamily\footnotesize,
  backgroundcolor=\color{codegray},
  frame=single,
  rulecolor=\color{rulegray},
  framerule=0.4pt,
  breaklines=true,
  breakatwhitespace=false,
  columns=fullflexible,
  keepspaces=true,
  showstringspaces=false,
  aboveskip=5pt,
  belowskip=5pt,
  xleftmargin=2pt,
  xrightmargin=2pt
}

\newcommand{\tip}[1]{\textbf{Tip.} #1}
\newcommand{\why}[1]{\textbf{Why.} #1}
\newcommand{\warn}[1]{\textbf{Watch out.} #1}

\title{\textbf{LaTeX Preamble \& Setup Quick Lookup}\\
\large Packages, settings, shortcuts, and practical tricks before \texttt{\textbackslash begin\{document\}}}
\author{For beginning LaTeX users in PhD coursework and research}
\date{}

\begin{document}
\maketitle
\vspace{-1.2em}

\begin{center}
\small
The focus of this handbook is the \textbf{setup area before writing content}. Optional lines are often shown commented with \texttt{\%}; remove the leading \texttt{\%} only when you need that feature.
\end{center}

\tableofcontents
\vspace{0.5em}
\hrule
\vspace{0.6em}

\section{How to use this guide if you know almost no LaTeX}

LaTeX has two broad parts:

\begin{enumerate}
\item the \textbf{preamble}: everything before \verb|\begin{document}|; this is where you load packages and define global rules;
\item the \textbf{document body}: everything between \verb|\begin{document}| and \verb|\end{document}|; this is where you write the paper.
\end{enumerate}

A useful mental model is:

\begin{center}
\texttt{document class} $\rightarrow$ \texttt{packages} $\rightarrow$ \texttt{global settings/macros}
$\rightarrow$ \texttt{\textbackslash begin\{document\}} $\rightarrow$ \texttt{paper content}.
\end{center}

\subsection{A tiny file that compiles}

\begin{lstlisting}
\documentclass[12pt]{article}
\usepackage[margin=1in]{geometry}
\usepackage{amsmath}

\title{My First Economics Note}
\author{Your Name}
\date{\today}

\begin{document}
\maketitle

This is ordinary text. Inline math looks like $u(c)$.

\begin{equation}
U = \sum_{t=0}^{T} \beta^t u(c_t).
\end{equation}

\end{document}
\end{lstlisting}

\textbf{What happens when you compile it?} LaTeX reads the setup first, then creates the title, text, and numbered equation. The source file is \texttt{.tex}; the usual finished output is a PDF.

\subsection{The three questions to ask before adding anything}

Before adding a package or setting, ask:

\begin{enumerate}[label=(\arabic*)]
\item \textbf{Do I actually need it?} Keep the preamble small while learning.
\item \textbf{Does my class/journal template already provide it?} Avoid fighting an official template.
\item \textbf{Is this global or local?} Global formatting belongs in the preamble; one-off formatting usually belongs in the body.
\end{enumerate}

\tip{For a new project, first make the smallest file compile. Add figures, bibliography, custom macros, and special formatting one module at a time.}

\section{Compiler / engine choice: pdfLaTeX, XeLaTeX, or LuaLaTeX}

For a beginner writing an English-language academic paper, \textbf{pdfLaTeX is usually the easiest default}. XeLaTeX and LuaLaTeX are especially useful when you need system fonts or multilingual text.

\begin{tabularx}{\textwidth}{>{\bfseries}p{0.18\textwidth} X X}
\toprule
Engine & Good default for & Typical font setup \\
\midrule
pdfLaTeX & Coursework, papers, journal-style PDFs & Traditional LaTeX font packages such as \texttt{newtxtext/newtxmath} \\
XeLaTeX & System fonts, Unicode-heavy documents, multilingual work & \texttt{fontspec} \\
LuaLaTeX & Similar benefits to XeLaTeX plus Lua-based features & \texttt{fontspec} \\
\bottomrule
\end{tabularx}

For modern pdfLaTeX installations, UTF-8 input is normally the default. You may still encounter older templates containing:

\begin{lstlisting}
% Mostly needed in older pdfLaTeX templates:
% \usepackage[utf8]{inputenc}
% \usepackage[T1]{fontenc}
\end{lstlisting}

For XeLaTeX/LuaLaTeX, a basic system-font setup looks like:

\begin{lstlisting}
% Compile with XeLaTeX or LuaLaTeX -- not pdfLaTeX
% \usepackage{fontspec}
% \setmainfont{Times New Roman}
\end{lstlisting}

\warn{Do not combine \texttt{fontspec} with pdfLaTeX. If a template already chooses an engine/font system, follow the template.}

\section{The 30-second starter preamble}

For most quantitative or technical homework, start here. It is intentionally small.

\begin{lstlisting}
\documentclass[12pt]{article}

% ----- Math -----
\usepackage{amsmath}     % align, cases, equation tools
\usepackage{amssymb}     % extra math symbols
\usepackage{amsfonts}    % \mathbb{R}, \mathbb{E}, etc.
\usepackage{mathtools}   % useful extensions to amsmath
\usepackage{bm}          % bold math: \bm{\beta}, \bm{x}

% ----- Page layout -----
\usepackage[margin=1in]{geometry}

% ----- Useful for economics tables -----
\usepackage{booktabs}

% ----- Hyperlinks / references -----
\usepackage[hidelinks]{hyperref}

\title{ECONS XXX -- Homework X}
\author{Your Name}
\date{\today}

\begin{document}
% Your content begins here.
\end{document}
\end{lstlisting}

\why{This covers clean math, standard margins, tables, and clickable PDF references without loading a large number of packages.}

\tip{Keep a master copy. Add optional modules only when needed instead of starting from a giant preamble every time.}

\section{How comments work in the setup area}

Anything after \verb|%| on a line is ignored by LaTeX. This makes comments perfect for keeping optional packages and reminders in your template.

\begin{lstlisting}
% OPTIONAL: remove the % when needed
% \usepackage{graphicx}      % figures
% \usepackage{float}         % [H] figure placement
% \usepackage{subcaption}    % subfigures
% \usepackage{setspace}      % line spacing
% \usepackage{fancyhdr}      % headers / footers
% \usepackage{siunitx}       % decimal alignment in tables
\end{lstlisting}

To activate a package, change

\begin{lstlisting}
% \usepackage{graphicx}
\end{lstlisting}

to

\begin{lstlisting}
\usepackage{graphicx}
\end{lstlisting}

\warn{Do not uncomment every package just because it is available. More packages can mean more conflicts and slower debugging.}

\section{Document class and global options}

The first line controls the overall document type.

\begin{lstlisting}
\documentclass[12pt]{article}

% Alternatives:
% \documentclass[11pt]{article}       % slightly more compact
% \documentclass[12pt]{report}        % chapters; longer reports
% \documentclass[12pt]{book}          % dissertation / book style

% Useful options:
% \documentclass[12pt,fleqn]{article} % display equations left-aligned
% \documentclass[12pt,leqno]{article} % equation numbers on the left
\end{lstlisting}

For most problem sets and research papers, \texttt{article} is usually the simplest choice.

\section{Page layout, margins, and paper size}

\subsection{Geometry}

\begin{lstlisting}
\usepackage{geometry}
\geometry{margin=1in}

% North America:
% \geometry{letterpaper, margin=1in}

% Many international settings:
% \geometry{a4paper, margin=2.5cm}

% Fine control:
% \geometry{
%   top=1in,
%   bottom=1in,
%   left=1.1in,
%   right=1.1in
% }
\end{lstlisting}

\tip{If an instructor or journal specifies margins, set them once here. Do not manually push text with repeated \texttt{\\hspace} or \texttt{\\vspace}.}

\subsection{Useful geometry patterns}

You can change only the dimension you need:

\begin{lstlisting}
% Wider left margin for binding / notes:
% \geometry{left=1.25in, right=1in, top=1in, bottom=1in}

% Add room between the header and body:
% \geometry{headsep=18pt}

% Include header/footer in the measured page area:
% \geometry{includeheadfoot, margin=1in}
\end{lstlisting}

\why{A page-layout package is more stable than scattering manual spaces throughout a paper. If the paper later changes from 1-inch to 1.25-inch margins, one preamble edit updates the entire document.}

\subsection{Paragraph style}

\begin{lstlisting}
% Default LaTeX: indented paragraphs, little vertical gap.

% Optional "block paragraph" style:
% \setlength{\parindent}{0pt}
% \setlength{\parskip}{6pt}
\end{lstlisting}

For formal papers, the default paragraph indentation is often safer unless a style guide says otherwise.

\section{Title, authors, affiliations, acknowledgments, and PDF metadata}

The simplest setup is:

\begin{lstlisting}
\title{The Economics of Something Important}
\author{Your Name}
\date{September 2026}

% In the body:
% \maketitle
\end{lstlisting}

A common paper-style pattern uses \verb|\thanks| for affiliations or acknowledgments:

\begin{lstlisting}
\title{Paper Title}
\author{
  First Author\thanks{Department of Economics, University A. Email: a@uni.edu.}
  \and
  Second Author\thanks{Department of Economics, University B. Email: b@uni.edu.}
}
\date{\today}
\end{lstlisting}

You can also set PDF metadata through \texttt{hyperref}:

\begin{lstlisting}
\usepackage[hidelinks]{hyperref}
\hypersetup{
  pdftitle={Paper Title},
  pdfauthor={Your Name},
  pdfsubject={Economics},
  pdfkeywords={labor economics, productivity, firms}
}
\end{lstlisting}

\tip{Do not manually type a title at the top of page 1 unless a template requires it. Defining \texttt{\textbackslash title}, \texttt{\textbackslash author}, and \texttt{\textbackslash date} keeps the source reusable.}

\section{Fonts and typography}

\subsection{Keep the default, or choose one family}

\begin{lstlisting}
% Default LaTeX font is already suitable for academic work.

% Times-like text + math:
% \usepackage{newtxtext}
% \usepackage{newtxmath}

% Palatino-like text + math:
% \usepackage{newpxtext}
% \usepackage{newpxmath}

% Latin Modern:
% \usepackage{lmodern}
\end{lstlisting}

\warn{Do not load several competing text/math font packages together. Pick one font system.}

\subsection{Microtype: an easy polish improvement}

\begin{lstlisting}
% Recommended for polished PDF typography:
% \usepackage{microtype}
\end{lstlisting}

It improves line breaking, spacing, and character protrusion with almost no effort.

\section{Line spacing and paragraph spacing}

\begin{lstlisting}
% \usepackage{setspace}

% \onehalfspacing     % 1.5 line spacing
% \doublespacing      % double spacing
% \singlespacing      % single spacing

% Local spacing only:
% \begin{spacing}{1.2}
% Text here...
% \end{spacing}
\end{lstlisting}

\tip{Use these settings in the preamble if your instructor requires 1.5 or double spacing. Avoid manually inserting blank lines everywhere.}

\section{Core math packages for technical writing}

A practical math block is:

\begin{lstlisting}
\usepackage{amsmath}     % align, split, cases, equation environments
\usepackage{amssymb}     % symbols such as \mathbb and many relations
\usepackage{amsfonts}    % math fonts
\usepackage{mathtools}   % extensions and fixes for amsmath
\usepackage{bm}          % bold vectors and Greek symbols

% Optional:
% \usepackage{cancel}    % show cancellation in derivations
% \usepackage{physics}   % convenient notation, but can redefine commands
\end{lstlisting}

Typical uses:

\begin{tabularx}{\textwidth}{>{\ttfamily}l X}
\toprule
Package & Why researchers often use it \\
\midrule
amsmath & Multi-line derivations, FOCs, systems, equation numbering \\
amssymb / amsfonts & Blackboard bold sets and many symbols \\
mathtools & Better delimiters, math extensions, small fixes \\
bm & Bold vectors such as \texttt{\textbackslash bm\{x\}} and \texttt{\textbackslash bm\{beta\}} \\
\bottomrule
\end{tabularx}

\subsection{What these packages let you write}

Once the math packages are loaded, ordinary quantitative notation becomes straightforward. For example:

\begin{lstlisting}
\[
\E[Y \mid X=x], \qquad
\bm{\beta} = (\beta_1,\ldots,\beta_K)', \qquad
A \in \mathbb{R}^{n\times k}.
\]
\end{lstlisting}

In the compiled PDF, the intended mathematical structure is visible as
\[
\mathbb{E}[Y \mid X=x], \qquad
\bm{\beta} = (\beta_1,\ldots,\beta_K)', \qquad
A \in \mathbb{R}^{n\times k}.
\]

Here \verb|\qquad| is simply a convenient large horizontal math space; it is not part of the mathematics itself.

\section{Create your own notation shortcuts in the preamble}

Repeated notation becomes easier to read and maintain if you define it once.

\begin{lstlisting}
% ---------- Probability and sets ----------
% \newcommand{\E}{\mathbb{E}}
% \newcommand{\Prob}{\mathbb{P}}
% \newcommand{\R}{\mathbb{R}}
% \newcommand{\N}{\mathbb{N}}

% ---------- Operators ----------
% \DeclareMathOperator{\Var}{Var}
% \DeclareMathOperator{\Cov}{Cov}
% \DeclareMathOperator{\Corr}{Corr}
% \DeclareMathOperator{\plim}{plim}
% \DeclareMathOperator*{\argmax}{arg\,max}
% \DeclareMathOperator*{\argmin}{arg\,min}

% ---------- Differential ----------
% \newcommand{\dd}{\,\mathrm{d}}

% ---------- Common bold notation ----------
% \newcommand{\bx}{\bm{x}}
% \newcommand{\by}{\bm{y}}
% \newcommand{\bX}{\bm{X}}
% \newcommand{\bbeta}{\bm{\beta}}
% \newcommand{\beps}{\bm{\varepsilon}}
\end{lstlisting}

After activation, you can write concise source such as

\begin{lstlisting}
% \E[Y \mid X], \Var(X), X \in \R
% \plim \hat{\beta}, \argmax_x U(x)
\end{lstlisting}

\why{If you later change notation, you edit one definition rather than hundreds of equations.}

\subsection{Commands with arguments}

A macro can accept input. This is useful for notation that has a repeated structure:

\begin{lstlisting}
% One argument:
% \newcommand{\Econd}[1]{\mathbb{E}\!\left[#1\right]}
%
% Body:
% \Econd{Y_i \mid X_i}
\end{lstlisting}

which produces the same style every time:
\[
\mathbb{E}\!\left[Y_i\mid X_i\right].
\]

With two arguments:

\begin{lstlisting}
% \newcommand{\deriv}[2]{\frac{\partial #1}{\partial #2}}
% Body: \deriv{U}{c}
\end{lstlisting}

\subsection{Paired delimiters with mathtools}

For absolute values and norms that can resize automatically:

\begin{lstlisting}
% Requires mathtools
% \DeclarePairedDelimiter{\abs}{\lvert}{\rvert}
% \DeclarePairedDelimiter{\norm}{\lVert}{\rVert}
%
% \abs{x}
% \abs*{\frac{x}{y}}     % starred form auto-sizes
% \norm*{\bm{\beta}}
\end{lstlisting}

\warn{Avoid redefining a very common built-in command unless you know exactly what it does. Choose descriptive macro names and keep all your custom definitions in one block.}

\section{Equation numbering and reference setup}

\begin{lstlisting}
% Number equations by section:
% \numberwithin{equation}{section}

% Example labels in the document body:
% \begin{equation}
% y_i = \beta_0 + \beta_1 x_i + \varepsilon_i
% \label{eq:regression}
% \end{equation}

% Refer to it with:
% Equation~\eqref{eq:regression}
\end{lstlisting}

Recommended label prefixes:

\begin{tabular}{ll@{\qquad}ll}
\toprule
\texttt{eq:} & equation & \texttt{fig:} & figure \\
\texttt{tab:} & table & \texttt{sec:} & section \\
\texttt{app:} & appendix & \texttt{thm:} & theorem \\
\bottomrule
\end{tabular}

\tip{Put \texttt{\\label\{...\}} immediately after \texttt{\\caption\{...\}} in figures and tables.}

\subsection{A complete equation-reference example}

\begin{lstlisting}
\begin{equation}
  y_i = \alpha + \beta x_i + \varepsilon_i,
  \label{eq:baseline}
\end{equation}

As shown in Equation~\eqref{eq:baseline}, ...
\end{lstlisting}

The label is a stable internal name. If the equation moves from (2) to (5), the reference updates automatically after recompilation.

For multi-line mathematics:

\begin{lstlisting}
\begin{align}
U(c_1,c_2) &= u(c_1) + \beta u(c_2), \\
\text{s.t.}\qquad c_1 + \frac{c_2}{1+r} &= y_1 + \frac{y_2}{1+r}.
\end{align}
\end{lstlisting}

\tip{Put \texttt{\&} immediately before the symbol you want aligned, commonly \texttt{=} or \texttt{\textbackslash le}. Use \texttt{\textbackslash notag} on a line that should not receive an equation number.}

\section{Hyperlinks and smart references}

\begin{lstlisting}
% Usually load near the end of the preamble:
\usepackage[hidelinks]{hyperref}

% Optional smart references -- load AFTER hyperref:
% \usepackage[nameinlink,noabbrev]{cleveref}
\end{lstlisting}

With \texttt{cleveref}, source can become:

\begin{lstlisting}
% \cref{fig:eventstudy}
% \cref{tab:mainreg}
% \cref{eq:FOC}
\end{lstlisting}

instead of manually typing ``Figure'', ``Table'', or ``Equation'' each time.

\subsection{The nonbreaking-space trick}

\begin{lstlisting}
Figure~\ref{fig:main}
Table~\ref{tab:main}
Section~\ref{sec:data}
Equation~\eqref{eq:model}
\end{lstlisting}

The \verb|~| prevents the label word and its number from splitting across lines.

\subsection{Manual vs. automatic reference names}

\begin{tabularx}{\textwidth}{>{\ttfamily}p{0.30\textwidth} X}
\toprule
Source & Typical output / idea \\
\midrule
Figure\~\textbackslash ref\{fig:main\} & You supply the word ``Figure'' yourself. \\
\textbackslash autoref\{fig:main\} & \texttt{hyperref} supplies the object name. \\
\textbackslash cref\{fig:main\} & \texttt{cleveref} supplies the object name and handles multiple labels well. \\
\bottomrule
\end{tabularx}

\tip{For a large paper with many tables, figures, equations, appendices, and propositions, \texttt{cleveref} can substantially reduce manual reference formatting.}

\section{Tables: a professional research-paper setup}

\subsection{Packages worth keeping available}

\begin{lstlisting}
\usepackage{booktabs}        % professional horizontal rules

% \usepackage{array}         % custom column types
% \usepackage{tabularx}      % tables fit to a target width
% \usepackage{longtable}     % tables spanning multiple pages
% \usepackage{multirow}      % cells spanning rows
% \usepackage{siunitx}       % decimal alignment / number formatting
% \usepackage{threeparttable}% table notes
\end{lstlisting}

\why{Regression and summary-statistics tables are common in quantitative research papers. \texttt{booktabs} gives clean \texttt{\\toprule}, \texttt{\\midrule}, and \texttt{\\bottomrule}.}

\subsection{Decimal alignment with siunitx}

\begin{lstlisting}
% \usepackage{siunitx}
%
% \begin{tabular}{l S[table-format=2.2]}
% \toprule
% Variable & {Mean} \\
% \midrule
% Income   & 12.34 \\
% Age      & 30.10 \\
% \bottomrule
% \end{tabular}
\end{lstlisting}

\tip{For professional academic tables, avoid vertical lines unless specifically required.}

\subsection{A complete regression-style table pattern}

This is the body code you are preparing the preamble to support:

\begin{lstlisting}
\begin{table}[htbp]
\centering
\caption{Summary Statistics}
\label{tab:summary}
\begin{threeparttable}
\begin{tabular}{l S[table-format=4.1] S[table-format=3.2] S[table-format=6.0]}
\toprule
Variable & {Mean} & {SD} & {N} \\
\midrule
Income   & 4250.5 & 811.32 & 12500 \\
Age      &   38.7 &  12.41 & 12500 \\
\bottomrule
\end{tabular}
\begin{tablenotes}[flushleft]\footnotesize
\item Notes: This is a toy example showing table structure.
\end{tablenotes}
\end{threeparttable}
\end{table}
\end{lstlisting}

The central pieces are \texttt{booktabs} for clean rules, \texttt{siunitx} for numeric alignment, and \texttt{threeparttable} for notes that belong to the table.

\begin{center}
\begin{threeparttable}
\caption*{Compiled miniature example}
\begin{tabular}{l S[table-format=4.1] S[table-format=3.2] S[table-format=5.0]}
\toprule
Variable & {Mean} & {SD} & {N} \\
\midrule
Income & 4250.5 & 811.32 & 12500 \\
Age & 38.7 & 12.41 & 12500 \\
\bottomrule
\end{tabular}
\begin{tablenotes}[flushleft]\footnotesize
\item Notes: Numbers align cleanly without vertical rules.
\end{tablenotes}
\end{threeparttable}
\end{center}

\subsection{Useful siunitx setup for regression tables}

\begin{lstlisting}
% In the preamble:
% \sisetup{
%   detect-weight=true,
%   detect-inline-weight=math,
%   table-number-alignment=center
% }
\end{lstlisting}

\warn{If your regression table is generated by Stata/R, do not manually retype hundreds of coefficients into LaTeX. Export a \texttt{.tex} fragment and include it with \texttt{\textbackslash input}.}

\section{Figures, plots, and subfigures}

\begin{lstlisting}
% \usepackage{graphicx}     % \includegraphics
% \usepackage{float}        % enables [H]
% \usepackage{caption}      % caption control
% \usepackage{subcaption}   % subfigures
\end{lstlisting}

A useful folder shortcut:

\begin{lstlisting}
% Requires graphicx
% \graphicspath{{figures/}}
\end{lstlisting}

Then you can write \verb|\includegraphics{eventstudy.pdf}| instead of the full folder path.

Common float placement letters:

\begin{tabularx}{\textwidth}{>{\ttfamily}l X}
\toprule
Option & Meaning \\
\midrule
h & approximately here \\
t & top of page \\
b & bottom of page \\
p & dedicated float page \\
htbp & let LaTeX choose among sensible locations \\
H & exactly here; requires \texttt{float} \\
\bottomrule
\end{tabularx}

\warn{Do not use \texttt{[H]} for every figure. \texttt{[htbp]} usually produces a better paper layout.}

\subsection{A complete figure pattern}

\begin{lstlisting}
\begin{figure}[htbp]
  \centering
  \includegraphics[width=0.75\linewidth]{figures/main_result.pdf}
  \caption{Main Result}
  \label{fig:main}
\end{figure}

See Figure~\ref{fig:main}.
\end{lstlisting}

\textbf{Why \texttt{\textbackslash linewidth}?} Inside ordinary text it is usually the available line width; inside a smaller container such as a \texttt{minipage}, it automatically adapts to that container. This is often safer than hard-coding inches.

Common sizing patterns:

\begin{lstlisting}
% Fill available local width:
% \includegraphics[width=\linewidth]{figure.pdf}

% Use 70% of available local width:
% \includegraphics[width=.70\linewidth]{figure.pdf}

% Limit height while preserving aspect ratio:
% \includegraphics[height=.55\textheight,keepaspectratio]{figure.pdf}

% Crop whitespace: trim = left bottom right top
% \includegraphics[width=.8\linewidth,trim=10 5 10 5,clip]{figure.pdf}
\end{lstlisting}

\subsection{Subfigures / panels}

\begin{lstlisting}
% Preamble:
% \usepackage{subcaption}
%
% Body:
% \begin{figure}[htbp]
% \centering
% \begin{subfigure}{.48\linewidth}
%   \includegraphics[width=\linewidth]{panel_a.pdf}
%   \caption{Panel A}
%   \label{fig:panel-a}
% \end{subfigure}\hfill
% \begin{subfigure}{.48\linewidth}
%   \includegraphics[width=\linewidth]{panel_b.pdf}
%   \caption{Panel B}
%   \label{fig:panel-b}
% \end{subfigure}
% \caption{Two Related Outcomes}
% \label{fig:two-panels}
% \end{figure}
\end{lstlisting}

\section{Control floating tables and figures}

If LaTeX moves a table too far away from the relevant section:

\begin{lstlisting}
% \usepackage[section]{placeins}
\end{lstlisting}

This inserts barriers at sections. Or use a manual barrier:

\begin{lstlisting}
% \FloatBarrier
\end{lstlisting}

\section{Bibliography and citation setup}

Choose \textbf{one} main workflow.

\subsection{Option A: biblatex + biber}

\begin{lstlisting}
% \usepackage{csquotes}
% \usepackage[
%   backend=biber,
%   style=authoryear,
%   natbib=true,
%   maxcitenames=2,
%   maxbibnames=99
% ]{biblatex}
%
% \addbibresource{references.bib}
\end{lstlisting}

Common body commands:

\begin{lstlisting}
% \textcite{angrist2009}     % Angrist and Pischke (2009)
% \parencite{angrist2009}    % (Angrist and Pischke, 2009)
% \printbibliography
\end{lstlisting}

\subsection{Option B: natbib + BibTeX}

\begin{lstlisting}
% \usepackage[round,authoryear]{natbib}
%
% Body:
% \citet{angrist2009}        % Angrist and Pischke (2009)
% \citep{angrist2009}        % (Angrist and Pischke, 2009)
%
% End of document:
% \bibliographystyle{apalike}
% \bibliography{references}
\end{lstlisting}

\warn{Do not independently configure both \texttt{biblatex} and \texttt{natbib} in the same document. Choose the workflow required by your class, coauthors, or journal.}

\subsection{What goes inside references.bib?}

A typical article entry looks like:

\begin{lstlisting}
@article{acemoglu2001,
  author  = {Daron Acemoglu and Joshua Angrist},
  title   = {Consequences of Employment Protection?},
  journal = {Journal of Political Economy},
  year    = {2001},
  volume  = {109},
  number  = {5},
  pages   = {915--957}
}
\end{lstlisting}

The key \texttt{acemoglu2001} is the identifier used by your citation command. For example:

\begin{lstlisting}
\citet{acemoglu2001}
\citep{acemoglu2001}
% or, with biblatex:
\textcite{acemoglu2001}
\end{lstlisting}

\subsection{Compile sequence: why the bibliography sometimes appears only after several runs}

With BibTeX, the conceptual sequence is:

\begin{lstlisting}
pdflatex main
bibtex main
pdflatex main
pdflatex main
\end{lstlisting}

With \texttt{biblatex} + Biber:

\begin{lstlisting}
pdflatex main
biber main
pdflatex main
pdflatex main
\end{lstlisting}

If you use \texttt{latexmk}, it can usually determine the required repeated passes automatically:

\begin{lstlisting}
latexmk -pdf main.tex
\end{lstlisting}

\tip{In Overleaf, the editor usually handles the repeated passes for you. Locally, \texttt{latexmk} is one of the easiest ways to avoid thinking about the exact compilation sequence.}

\section{Theory papers: theorem, proposition, assumption}

\begin{lstlisting}
% \usepackage{amsthm}
%
% \newtheorem{theorem}{Theorem}
% \newtheorem{proposition}{Proposition}
% \newtheorem{lemma}{Lemma}
% \newtheorem{corollary}{Corollary}
%
% \theoremstyle{definition}
% \newtheorem{definition}{Definition}
% \newtheorem{assumption}{Assumption}
%
% \theoremstyle{remark}
% \newtheorem{remark}{Remark}
\end{lstlisting}

Number by section if desired:

\begin{lstlisting}
% \newtheorem{theorem}{Theorem}[section]
\end{lstlisting}

A body example:

\begin{lstlisting}
\begin{proposition}\label{prop:monotonicity}
If $u'(c)>0$, utility is increasing in consumption.
\end{proposition}

\begin{proof}
The result follows directly from the sign of the derivative.
\end{proof}
\end{lstlisting}

The point of the preamble declarations is that propositions, assumptions, lemmas, and proofs all receive consistent formatting and numbering throughout the paper.

\section{Headers, footers, and page numbering}

\begin{lstlisting}
% \usepackage{fancyhdr}
%
% \pagestyle{fancy}
% \fancyhf{}                 % clear existing header/footer fields
% \lhead{ECONS 5XX}
% \chead{Homework 3}
% \rhead{Your Name}
% \cfoot{\thepage}
%
% \renewcommand{\headrulewidth}{0pt}
% \renewcommand{\footrulewidth}{0pt}
\end{lstlisting}

Simpler alternatives:

\begin{lstlisting}
% \pagestyle{plain}          % page number only
% \pagestyle{empty}          % no page numbers
\end{lstlisting}

\section{Section-heading control}

\begin{lstlisting}
% \usepackage{titlesec}
%
% Compact spacing:
% \titlespacing*{\section}{0pt}{1.5ex}{0.8ex}
%
% Custom style:
% \titleformat{\section}
%   {\large\bfseries}
%   {\thesection}{1em}{}
\end{lstlisting}

\warn{Journal templates often control headings themselves. If you are using an official class file, avoid overriding its style unless permitted.}

\section{Lists and compact outlines}

\begin{lstlisting}
% \usepackage{enumitem}
%
% Compact list defaults:
% \setlist[itemize]{noitemsep, topsep=2pt}
% \setlist[enumerate]{noitemsep, topsep=2pt}
%
% Lettered list:
% \begin{enumerate}[label=(\alph*)]
%   \item First item
% \end{enumerate}
\end{lstlisting}

Useful for assumptions, model steps, homework subparts, and referee-response style documents.

\section{Appendix and landscape pages}

\subsection{Appendix}

No package is required for the basic approach:

\begin{lstlisting}
% In the document body:
% \appendix
% \section{Additional Results}
\end{lstlisting}

Optional enhancements:

\begin{lstlisting}
% \usepackage[toc,page]{appendix}
\end{lstlisting}

\subsection{Wide regression tables}

\begin{lstlisting}
% \usepackage{pdflscape}
%
% \begin{landscape}
%   ... very wide table ...
% \end{landscape}
\end{lstlisting}

This is better than manually rotating every cell.

\section{Stata, R, Python, and external regression tables}

For code snippets:

\begin{lstlisting}
% \usepackage{listings}
\end{lstlisting}

For tables exported by Stata/R/Python, keep generated tables in separate files:

\begin{lstlisting}
% Main paper:
% \input{tables/regression_main.tex}
% \input{tables/summary_statistics.tex}
\end{lstlisting}

Recommended project structure:

\begin{lstlisting}
project/
  main.tex
  references.bib
  sections/
    introduction.tex
    model.tex
    data.tex
    results.tex
  figures/
    eventstudy.pdf
    trends.pdf
  tables/
    summary_statistics.tex
    regression_main.tex
\end{lstlisting}

Then assemble the paper cleanly:

\begin{lstlisting}
% \input{sections/introduction}
% \input{sections/model}
% \input{sections/data}
% \input{sections/results}
\end{lstlisting}

\why{Separating tables, figures, and sections keeps a research project much easier to maintain and reproduce.}

\section{Small spacing and alignment commands worth recognizing}

These are mostly body-level commands, but they are worth knowing because they often appear in examples and custom macros in the preamble.

\begin{tabularx}{\textwidth}{>{\ttfamily}p{0.18\textwidth} X X}
\toprule
Command & What it does & Typical use \\
\midrule
\textbackslash, & tiny math space & $f(x)\,\mathrm{d}x$ \\
\textbackslash; & larger math space & separate nearby math objects \\
\textbackslash! & negative math space & tighten a wide construction \\
\textbackslash quad & large fixed math space & separate equation and condition \\
\textbackslash qquad & about twice \texttt{\textbackslash quad} & $x=1\qquad i=1,\ldots,n$ \\
\textasciitilde & nonbreaking text space & \texttt{Table\textasciitilde\textbackslash ref\{...\}} \\
\textbackslash hfill & stretchable horizontal space & push material to opposite sides \\
\textbackslash noindent & suppress one paragraph indent & special front matter / notes \\
\& & alignment point in tables/\texttt{align} & align equations at equals signs \\
\textbackslash\textbackslash & end row/line in structured environments & tables and \texttt{align} \\
\bottomrule
\end{tabularx}

A classic quantitative-writing use of \verb|\qquad| is:

\begin{lstlisting}
\[
y_i = x_i'\beta + \varepsilon_i,
\qquad i=1,\ldots,n.
\]
\end{lstlisting}

which appears as
\[
y_i = x_i'\beta + \varepsilon_i,
\qquad i=1,\ldots,n.
\]

\warn{Spacing commands should improve readability, not manually ``draw'' the page. Use semantic structures such as \texttt{align}, \texttt{tabular}, \texttt{minipage}, and page-layout packages before resorting to repeated manual spacing.}

\section{Draft-mode tricks that save time}

\subsection{Show overfull lines}

\begin{lstlisting}
% Draw a black marker where text exceeds the margin:
% \overfullrule=5pt
\end{lstlisting}

Remove it before final submission.

\subsection{Show label names}

\begin{lstlisting}
% \usepackage{showkeys}
\end{lstlisting}

This displays label keys such as \texttt{eq:ols} or \texttt{fig:eventstudy} while drafting.

\subsection{TODO notes}

\begin{lstlisting}
% \usepackage[colorinlistoftodos]{todonotes}
%
% Body examples:
% \todo{Check this result.}
% \todo[inline]{Rewrite this argument.}
%
% Final version switch:
% \usepackage[disable]{todonotes}
\end{lstlisting}

\subsection{One draft/final switch}

\begin{lstlisting}
% \newif\ifdraft
% \drafttrue
% % \draftfalse
%
% \ifdraft
%   % draft-only packages/settings here
% \fi
\end{lstlisting}

\tip{A single switch is much easier than manually deleting draft annotations throughout a long paper.}

\section{Common compilation and debugging tricks}

\begin{itemize}
\item If references show \texttt{??}, compile again. Cross-references generally need multiple passes.
\item If a bibliography is missing, make sure you are using the correct BibTeX/Biber workflow and compile sequence.
\item A missing closing brace \verb|}| may cause an error many lines after the true mistake.
\item If a figure is missing, first check the filename, extension, relative path, and whether \texttt{graphicx} is loaded.
\item If a command is undefined, check which package provides it.
\item If a package conflict appears, simplify the preamble and reintroduce optional packages one at a time.
\end{itemize}

Useful source-display trick:

\begin{lstlisting}
% Show raw LaTeX inline:
% \verb|\frac{a}{b}|
%
% Show a raw block:
% \begin{verbatim}
% \frac{a}{b}
% \end{verbatim}
\end{lstlisting}

\subsection{Common error messages: what to inspect first}

\begin{tabularx}{\textwidth}{>{\ttfamily}p{0.30\textwidth} X}
\toprule
Message / symptom & First checks \\
\midrule
Undefined control sequence & Misspelling? Missing package? Command used with wrong engine? \\
Missing \$ inserted & Math command outside math mode, or unmatched \texttt{\$}. \\
Extra alignment tab has been changed to \textbackslash cr & Too many \texttt{\&} symbols for the table/alignment structure. \\
Missing \} inserted & Unbalanced braces; inspect a few lines before the reported location. \\
Runaway argument? & Usually a missing closing brace or environment end. \\
Reference shows ?? & Compile again; confirm the label exists and is unique. \\
Citation shows ? & Check bibliography key and BibTeX/Biber step. \\
Overfull \textbackslash hbox & A line cannot break; inspect long URLs, code, math, or fixed-width content. \\
Float(s) lost / figures far away & Check float environments, placement options, barriers, and oversized figures. \\
\bottomrule
\end{tabularx}

\subsection{A practical local compile command}

If you have a TeX distribution installed locally:

\begin{lstlisting}
latexmk -pdf main.tex
\end{lstlisting}

To clean auxiliary files later:

\begin{lstlisting}
latexmk -c
\end{lstlisting}

\section{Package order: a practical rule}

A generally safe order is:

\begin{enumerate}
\item language / encoding,
\item fonts,
\item mathematics,
\item geometry and spacing,
\item tables and figures,
\item bibliography,
\item colors,
\item \texttt{hyperref},
\item \texttt{cleveref}.
\end{enumerate}

Example skeleton:

\begin{lstlisting}
\documentclass[12pt]{article}

% 1. Language / font
% \usepackage[english]{babel}
% \usepackage{newtxtext,newtxmath}

% 2. Math
\usepackage{amsmath,amssymb,amsfonts,mathtools,bm}

% 3. Layout
\usepackage[margin=1in]{geometry}
% \usepackage{setspace}

% 4. Tables / figures
\usepackage{booktabs}
% \usepackage{graphicx}
% \usepackage{siunitx}

% 5. Bibliography
% ... choose biblatex OR natbib ...

% 6. Links near the end
\usepackage[hidelinks]{hyperref}
% \usepackage[nameinlink,noabbrev]{cleveref}
\end{lstlisting}

\warn{Package manuals and official journal templates override this general rule. If a class file gives instructions, follow them.}

\section{Common beginner mistakes in the preamble}

\begin{tabularx}{\textwidth}{>{\bfseries}p{0.30\textwidth} X}
\toprule
Mistake & Better approach \\
\midrule
Loading every package & Start small; activate modules only when needed. \\
Mixing bibliography systems & Choose \texttt{biblatex/biber} or \texttt{natbib/BibTeX}. \\
Mixing several font packages & Choose one coherent text/math font setup. \\
Loading \texttt{cleveref} before \texttt{hyperref} & Usually load \texttt{hyperref} first, then \texttt{cleveref}. \\
Forcing all floats with \texttt{[H]} & Prefer \texttt{[htbp]} and let LaTeX manage placement. \\
Manual formatting everywhere & Put global rules in the preamble once. \\
Hard-coding repeated notation & Define macros with \texttt{\textbackslash newcommand} or \texttt{\textbackslash DeclareMathOperator}. \\
One giant \texttt{main.tex} & Split sections and exported tables with \texttt{\textbackslash input}. \\
Using \texttt{\textbackslash\textbackslash} for normal paragraph spacing & End a paragraph with a blank line; reserve \texttt{\textbackslash\textbackslash} for environments where a line/row break is meaningful. \\
Typing figure/table numbers manually & Use \texttt{\textbackslash label} and \texttt{\textbackslash ref} / \texttt{\textbackslash cref}. \\
Using \texttt{\$\$ ... \$\$} & Prefer \texttt{\textbackslash[ ... \textbackslash]} for unnumbered display math or \texttt{equation/align} for numbered math. \\
Hard-coding figure widths in inches everywhere & Prefer proportions such as \texttt{width=.8\textbackslash linewidth}. \\
Putting journal-specific formatting into every project & Keep a simple personal template and let official journal/class files control final styling. \\
\bottomrule
\end{tabularx}

\section{Ready-to-copy quantitative homework preamble}

This is a good default when you want something simple but useful.

\begin{lstlisting}
\documentclass[12pt]{article}

% ---------- Math ----------
\usepackage{amsmath}
\usepackage{amssymb}
\usepackage{amsfonts}
\usepackage{mathtools}
\usepackage{bm}

% ---------- Layout ----------
\usepackage[margin=1in]{geometry}

% ---------- Tables ----------
\usepackage{booktabs}

% ---------- Links ----------
\usepackage[hidelinks]{hyperref}

% ---------- Common quantitative notation ----------
\newcommand{\E}{\mathbb{E}}
\newcommand{\Prob}{\mathbb{P}}
\newcommand{\R}{\mathbb{R}}
\DeclareMathOperator{\Var}{Var}
\DeclareMathOperator{\Cov}{Cov}
\DeclareMathOperator{\plim}{plim}
\DeclareMathOperator*{\argmax}{arg\,max}
\DeclareMathOperator*{\argmin}{arg\,min}

% ---------- Optional modules ----------
% \usepackage{graphicx}      % figures
% \usepackage{float}         % [H]
% \usepackage{siunitx}       % decimal-aligned tables
% \usepackage{setspace}      % line spacing
% \usepackage{fancyhdr}      % headers / footers

\title{ECONS XXX -- Homework X}
\author{Your Name}
\date{\today}

\begin{document}
\maketitle

% Start your homework here.

\end{document}
\end{lstlisting}

\section{Ready-to-copy research-paper preamble}

This adds the pieces commonly useful for an empirical paper while remaining manageable.

\begin{lstlisting}
\documentclass[12pt]{article}

% ---------- Typography / layout ----------
\usepackage[margin=1in]{geometry}
\usepackage{microtype}

% ---------- Mathematics ----------
\usepackage{amsmath}
\usepackage{amssymb}
\usepackage{amsfonts}
\usepackage{mathtools}
\usepackage{bm}

% ---------- Figures / tables ----------
\usepackage{graphicx}
\usepackage{booktabs}
\usepackage{siunitx}
\usepackage{threeparttable}

% ---------- Optional paper tools ----------
% \usepackage{setspace}
% \onehalfspacing
% \usepackage{pdflscape}
% \usepackage[section]{placeins}

% ---------- Bibliography: choose ONE system ----------
% \usepackage[round,authoryear]{natbib}
%
% OR:
% \usepackage{csquotes}
% \usepackage[backend=biber,style=authoryear]{biblatex}
% \addbibresource{references.bib}

% ---------- Links / references ----------
\usepackage[hidelinks]{hyperref}
% \usepackage[nameinlink,noabbrev]{cleveref}

% ---------- Example quantitative notation ----------
\newcommand{\E}{\mathbb{E}}
\newcommand{\Prob}{\mathbb{P}}
\newcommand{\R}{\mathbb{R}}
\DeclareMathOperator{\Var}{Var}
\DeclareMathOperator{\Cov}{Cov}
\DeclareMathOperator{\Corr}{Corr}
\DeclareMathOperator{\plim}{plim}
\DeclareMathOperator*{\argmax}{arg\,max}
\DeclareMathOperator*{\argmin}{arg\,min}

% ---------- Project folders ----------
\graphicspath{{figures/}}

\title{Paper Title}
\author{Your Name}
\date{\today}

\begin{document}
\maketitle

% \input{sections/introduction}
% \input{sections/model}
% \input{sections/data}
% \input{sections/results}

\end{document}
\end{lstlisting}

\section{Ready-to-copy theory-paper preamble}

For a theory note or model-heavy field paper, theorem environments and notation consistency become more important than regression-table machinery.

\begin{lstlisting}
\documentclass[12pt]{article}

% ---------- Layout / typography ----------
\usepackage[margin=1in]{geometry}
\usepackage{microtype}

% ---------- Math ----------
\usepackage{amsmath,amssymb,amsfonts,mathtools,bm}
\usepackage{amsthm}

% ---------- Theorem environments ----------
\newtheorem{theorem}{Theorem}[section]
\newtheorem{proposition}[theorem]{Proposition}
\newtheorem{lemma}[theorem]{Lemma}
\newtheorem{corollary}[theorem]{Corollary}
\theoremstyle{definition}
\newtheorem{assumption}[theorem]{Assumption}
\newtheorem{definition}[theorem]{Definition}
\theoremstyle{remark}
\newtheorem{remark}[theorem]{Remark}

% ---------- Notation ----------
\newcommand{\E}{\mathbb{E}}
\newcommand{\R}{\mathbb{R}}
\DeclareMathOperator*{\argmax}{arg\,max}
\DeclareMathOperator*{\argmin}{arg\,min}
\DeclarePairedDelimiter{\abs}{\lvert}{\rvert}
\DeclarePairedDelimiter{\norm}{\lVert}{\rVert}

% ---------- References ----------
\usepackage[hidelinks]{hyperref}
% \usepackage[nameinlink,noabbrev]{cleveref}

\title{Theory Paper Title}
\author{Your Name}
\date{\today}

\begin{document}
\maketitle

% \input{sections/model}
% \input{sections/equilibrium}
% \input{sections/proofs}

\end{document}
\end{lstlisting}

\section{Module checklist: add only what the project needs}

\begin{tabularx}{\textwidth}{>{\bfseries}p{0.24\textwidth} p{0.28\textwidth} X}
\toprule
Need & Package / setup & Typical reason \\
\midrule
Basic math & \texttt{amsmath, amssymb, mathtools} & Equations, alignments, symbols \\
Bold Greek/vector math & \texttt{bm} & $\bm{\beta}$, $\bm{x}$ \\
Professional tables & \texttt{booktabs} & Clean regression/summary tables \\
Decimal alignment & \texttt{siunitx} & Numeric columns \\
Table notes & \texttt{threeparttable} & Notes under empirical tables \\
Figures & \texttt{graphicx} & Import PDF/PNG/JPG graphics \\
Figure panels & \texttt{subcaption} & Panel A / Panel B \\
Float barriers & \texttt{placeins} & Keep floats near a section \\
Landscape pages & \texttt{pdflscape} & Very wide tables \\
Theorems/proofs & \texttt{amsthm} & Theory papers \\
Line spacing & \texttt{setspace} & Referee/course requirements \\
Bibliography & \texttt{natbib} or \texttt{biblatex} & Citations and references \\
Smart refs & \texttt{cleveref} after \texttt{hyperref} & Automatic ``Figure/Table/Equation'' names \\
Draft TODOs & \texttt{todonotes} & Visible drafting reminders \\
\bottomrule
\end{tabularx}

\section{My recommended master-template strategy}

Keep one file with three layers:

\begin{enumerate}
\item \textbf{Always on:} core math, geometry, and one table package.
\item \textbf{Commented modules:} figures, bibliography, theorem tools, landscape pages, code, TODOs, special spacing.
\item \textbf{Your own macros:} notation that you repeatedly use in your field.
\end{enumerate}

A compact master block might look like this:

\begin{lstlisting}
% ==================== ALWAYS ON ====================
\usepackage{amsmath,amssymb,amsfonts,mathtools,bm}
\usepackage[margin=1in]{geometry}
\usepackage{booktabs}
\usepackage[hidelinks]{hyperref}

% ==================== OPTIONAL =====================
% \usepackage{graphicx}       % figures
% \usepackage{float}          % [H]
% \usepackage{subcaption}     % subfigures
% \usepackage{siunitx}        % aligned numbers
% \usepackage{threeparttable} % table notes
% \usepackage{setspace}       % 1.5 / double spacing
% \usepackage{fancyhdr}       % headers / footers
% \usepackage{amsthm}         % theorem environments
% \usepackage{pdflscape}      % landscape pages

% ==================== ECON MACROS ==================
% \newcommand{\E}{\mathbb{E}}
% \newcommand{\Prob}{\mathbb{P}}
% \DeclareMathOperator{\Var}{Var}
% \DeclareMathOperator{\Cov}{Cov}
% \DeclareMathOperator{\plim}{plim}
% \DeclareMathOperator*{\argmax}{arg\,max}
\end{lstlisting}

\end{document}
